\uuid{KiEI}
\titre{Extremum}
\auteur{bodin ; hanani ; mzari}
\organisation{exo7}
\licence{CC BY-NC-SA 4.0}
\datecreate{2020-01-20}
\niveau{L1}
\module{Analyse}
\chapitre{Dérivabilité des fonctions réelles}
\sousChapitre{Autre}
\difficulte{3}
\theme{Dérivées, extremum}
\exo7qcmid{lille-1.derivees.027}

\contenu{
\begin{question}
\qclassification{124.00}
\qtags{motcle=extremum}
Soit $n\geq 2$ un entier et $\displaystyle f(x)=\frac{1+x^n}{(1+x)^n}$. Quelles sont les bonnes réponses ?
\begin{answers}
    \good{$\displaystyle f'(x)=\frac{n(x^{n-1}-1)}{(1+x)^{n+1}}$ et $f$ admet un minimum en $1$.}
    \bad{$f'(1)\neq 0$ et donc $f$ n'admet pas d'extremum en $1$.}
    \bad{Le théorème de Rolle s'applique à $f$ sur $[-1,1]$ car $f(-1)=f(1)$.}
    \good{$\forall x\geq 0,\; (1+x)^n\leq 2^{n-1}(1+x^n)$.}
\end{answers}
\begin{explanations}
La dérivée de $f(x)$ est $\displaystyle f'(x)=\frac{n(x^{n-1}-1)}{(1+x)^{n+1}}$ et $f$ admet bien un minimum en $1$ (dresser le tableau de variations de $f$). En particulier,
$$\forall x\geq 0,\; f(1)\leq f(x)\Leftrightarrow (1+x)^n\leq 2^{n-1}(1+x^n).$$
\end{explanations}
\end{question}
}
